\section{数据集}
\label{ap:dataset_detail}


\subsection{\ourbenchmark}


\subsubsection{详细统计}

\begin{table}[h]
\centering
\small
\setlength{\tabcolsep}{6pt}
\begin{tabular}{c cc cc cc cc cc}
\toprule
\multirow{2}{*}{\textbf{取值}}
& \multicolumn{2}{c}{\textbf{深度}}
& \multicolumn{2}{c}{\textbf{时域}}
& \multicolumn{2}{c}{\textbf{广度}}
& \multicolumn{2}{c}{\textbf{并行性}}
& \multicolumn{2}{c}{\textbf{鲁棒性}} \\
\cmidrule(lr){2-3}
\cmidrule(lr){4-5}
\cmidrule(lr){6-7}
\cmidrule(lr){8-9}
\cmidrule(lr){10-11}
& 训练 & 测试
& 训练 & 测试
& 训练 & 测试
& 训练 & 测试
& 训练 & 测试 \\
\midrule

2  & \multirow{4}{*}{3993} & 199
   &                    & 167
   &                    & 200
   &                    & 167
   & \multirow{3}{*}{3000} & 200 \\

4  &                      & 200
   & \multirow{2}{*}{2174} & 150
   & \multirow{2}{*}{2000} & 200
   & \multirow{2}{*}{1807} & 150
   &                      & 200 \\

6  &                      & 200
   &                      & 126
   &                      & 192
   &                      & 126
   &                      & 200 \\

8  &                      & 200
   & ---                  & 124
   & ---                  & 84
   & ---                  & 124
   & ---                  & --- \\

10 & ---                  & 199
   & ---                  & ---
   & ---                  & ---
   & ---                  & ---
   & ---                  & --- \\

12 & ---                  & 197
   & ---                  & ---
   & ---                  & ---
   & ---                  & ---
   & ---                  & --- \\

\bottomrule
\end{tabular}
\caption{训练集/测试集取值，其中训练集单元格按纵向合并。}
\label{tab:mas_axes_multirow}
\end{table}


\subsubsection{示例}
\label{ap:masbench_example}

我们为每个维度给出一个取值为 4 的具体示例。

\textbf{深度}

\begin{systemprompt}[title={问题}]
每份 Aldi 桃罐头的数量等于每份橄榄罐头的高粱、每份 Marc's 橄榄罐头、每份桃罐头的黑麦和每份桃罐头的配料数量之和。每份 Marc's 鱼罐头的数量等于每份橄榄罐头配料数量的 15 倍。每份 Super Saver 汤罐头的数量等于每份桃罐头的黑麦数量。每份 Aldi 鱼罐头的数量等于 17。每份桃罐头的黑麦数量等于 11。每份汤罐头的高粱数量比每份 Super Saver 汤罐头多 19。每份桃罐头的藜麦数量等于 18。每份橄榄罐头的高粱数量等于 21。每份 Super Saver 桃罐头的数量是每份桃罐头黑麦数量的 2 倍。每份 Marc's 橄榄罐头的数量是每份橄榄罐头高粱数量的 3 倍。Marc's 有多少产品？
\end{systemprompt}
% \begin{userprompt}[title={Solution}]
%  Define Canned Olives's Sorghum as w; so w = 21. Define Canned Olives's Ingredient as h; so h = w = 21. Define Marc's's Canned Fish as y; so y = 15 * h = 15 * 21 = 16. Define Marc's's Canned Olives as T; so T = 3 * w = 3 * 21 = 17. Define Marc's's Product as i; so i = T + y = 17 + 16 = 10.
% \end{userprompt}
\begin{userprompt}[title={答案}]
10
\end{userprompt}

\textbf{时域}


\begin{systemprompt}[title={问题}]


问题 1：每个温带阔叶林灰熊的数量等于每个原始森林孟加拉虎与每个山地森林树懒数量之差。每个温带阔叶林树懒的数量比每个山地森林树懒多 10。每个山地森林树懒的数量等于 4。每个树懒的呼吸道黏膜数量比每个山地森林树懒多 18。每个树懒的鼻腔数量是每个温带阔叶林树懒数量的 10 倍。每个灰熊的呼吸道黏膜数量比每个温带阔叶林树懒多 20。温带阔叶林树懒的值是多少？\\

问题 2：每个温带阔叶林灰熊的数量等于每个原始森林孟加拉虎与每个山地森林树懒数量之差。每个树懒的口咽数量比每个树懒的呼吸道黏膜多 15。每个 [answer1] 的数量比每个山地森林树懒多 10。每个山地森林树懒的数量等于 4。每个树懒的呼吸道黏膜数量比每个山地森林树懒多 18。每个温带阔叶林孟加拉虎的数量等于每个树懒的器官与每个树懒的呼吸道黏膜数量之差。树懒呼吸道黏膜的值是多少？\\

问题 3：每个温带阔叶林灰熊的数量等于每个原始森林孟加拉虎与每个山地森林树懒数量之差。每个 [answer1] 的数量比每个山地森林树懒多 10。每个山地森林树懒的数量等于 4。每个 [answer2] 的数量比每个山地森林树懒多 18。每个孟加拉虎的呼吸道黏膜数量是每个灰熊呼吸道黏膜、每个树懒鼻腔和每个温带阔叶林孟加拉虎数量之和的 12 倍。每个灰熊的口咽数量等于每个灰熊的呼吸道黏膜数量。每个原始森林孟加拉虎的数量等于每个灰熊的口咽数量。每个树懒的鼻腔数量是每个 [answer1] 的 10 倍。每个灰熊的呼吸道黏膜数量比每个 [answer1] 多 20。使用上一次计算的结果 [answer1]，令 [variable3] = [answer1]。原始森林孟加拉虎是多少？\\

问题 4：每个温带阔叶林灰熊的数量等于每个 [answer3] 与每个山地森林树懒数量之差。每个树懒的口咽数量比每个 [answer2] 多 15。每个 [answer1] 的数量比每个山地森林树懒多 10。每个山地森林树懒的数量等于 4。每个 [answer2] 的数量比每个山地森林树懒多 18。每个孟加拉虎的口咽数量等于每个温带阔叶林生物与每个温带阔叶林孟加拉虎数量之和。每个孟加拉虎的呼吸道黏膜数量是每个灰熊呼吸道黏膜、每个树懒鼻腔和每个温带阔叶林孟加拉虎数量之和的 12 倍。每个温带阔叶林孟加拉虎的数量等于每个树懒器官与每个 [answer2] 数量之差。使用上一次计算的结果 [answer2]，令 [variable4] = K + [answer2]。孟加拉虎的口咽是多少？\\

注：在本组问题中：

- [variablek] 表示求解问题 k 所需的计算变量。

- [answerk] 表示问题 k 的答案。

逐步解决所有问题，并按以下格式给出全部答案：

\#\#\# 最终答案

问题 1：\boxed{[answer1]}

问题 2：\boxed{[answer2]}

问题 3：\boxed{[answer3]}

问题 4：\boxed{[answer4]}

\end{systemprompt}
% \begin{userprompt}[title={Solution}]
% Sub-problem 1: What is the value of Temperate Broadleaf Forest's Sloth?

% Sub-problem 2: What is the value of Sloth's Respiratory Mucosa?

% Sub-problem 3: Using the result [answer1] from the previous calculation, [variable3] = [answer1]. What is Old-growth Forest's Bengal Tiger?

% Sub-problem 4: Using the result [answer2] from the previous calculation, [variable4] = K + [answer2]. What is Bengal Tiger's Oropharynx?

% \end{userprompt}
\begin{userprompt}[title={答案}]
问题 1：\boxed{14}

问题 2：\boxed{22}

问题 3：\boxed{11}

问题 4：\boxed{7}
\end{userprompt}



\textbf{广度}


\begin{systemprompt}[title={问题}]
每个女妖心脏的数量是曼谷伦披尼公园生物数量与曼谷伦披尼公园水螅数量之差的 8 倍。每个曼谷伦披尼公园水螅的数量等于 0。每个曼谷伦披尼公园女妖的数量等于 10。每个腔静脉角质形成细胞的数量是每个曼谷伦披尼公园女妖数量的 17 倍。每个新加坡滨海湾花园水螅的数量等于每个曼谷伦披尼公园水螅与每个曼谷伦披尼公园器官数量之和。每个腔静脉肌腱细胞的数量比每个曼谷伦披尼公园女妖多 15。每个水螅的腔静脉数量是每个曼谷伦披尼公园水螅数量的 21 倍。水螅有多少个器官？
 \end{systemprompt}
% \begin{userprompt}[title={Solution}]
%  Define Lumpini Park in Bangkok's Hydra as r; so r = 0. Define Hydra's Vena Cava as O; so O = 21 * r = 21 * 0 = 0. Define Hydra's Organs as E; so E = O = 0.
%  \end{userprompt}
\begin{userprompt}[title={答案}]
0
\end{userprompt}


\textbf{并行性}


\begin{systemprompt}[title={问题}]
问题 1：每个华盛顿特区昆虫馆鲨鱼缸的数量等于每个芝加哥昆虫馆海藻森林缸的数量。每个西雅图昆虫馆海藻森林缸的数量等于 16。每个芝加哥昆虫馆海藻森林缸的数量是每个西雅图昆虫馆海藻森林缸数量的 22 倍。芝加哥昆虫馆海藻森林缸的值是多少？\\

问题 2：每个西雅图昆虫馆水族箱的数量等于每个芝加哥昆虫馆鲨鱼缸的数量。每个芝加哥昆虫馆水族箱的数量比每个西雅图昆虫馆水族箱多 0。每个华盛顿特区昆虫馆水族箱的数量等于每个芝加哥昆虫馆鲨鱼缸的数量。每个芝加哥昆虫馆鲨鱼缸的数量等于 17。西雅图昆虫馆水族箱的值是多少？\\

问题 3：每个西雅图昆虫馆水族箱的数量等于每个芝加哥昆虫馆鲨鱼缸的数量。每个华盛顿特区昆虫馆水族箱的数量等于每个芝加哥昆虫馆鲨鱼缸的数量。每个芝加哥昆虫馆鲨鱼缸的数量等于 17。华盛顿特区昆虫馆水族箱的值是多少？\\

问题 4：每个华盛顿特区昆虫馆海藻森林缸的数量等于每个芝加哥昆虫馆围场的数量。华盛顿特区昆虫馆海藻森林缸的值是多少？\\

注：在本组问题中：

- 每个问题都相互独立，可以并行求解。

逐步解决所有问题，并按以下格式给出全部答案：

\#\#\# 最终答案

问题 1：\boxed{[answer1]}

问题 2：\boxed{[answer2]}

问题 3：\boxed{[answer3]}

问题 4：\boxed{[answer4]}
\end{systemprompt}
% \begin{userprompt}[title={Solution}]
% Sub-problem 1: What is the value of Insectarium of Chicago's Kelp Forest Tank?

% Sub-problem 2: What is the value of Insectarium of Seattle's Aquarium?

% Sub-problem 3: What is the value of Insectarium of Washington DC's Aquarium?

% Sub-problem 4: What is the value of Insectarium of Washington DC's Kelp Forest Tank?

% \end{userprompt}
\begin{userprompt}[title={答案}]
问题 1：\boxed{7}

问题 2：\boxed{17}

问题 3：\boxed{17}

问题 4：\boxed{18}
\end{userprompt}

\textbf{鲁棒性}



\begin{systemprompt}[title={问题}]
仔细阅读以下段落。\\

每条金枪鱼的肱骨数量是 Tracy 鸟舍袋鼠步道、金枪鱼桡骨与枪鱼尺骨数量之和的 15 倍。每个熊猫展区的枪鱼数量比 Tracy 鸟舍围场数量与骆驼院动物数量之差多 13。每个袋鼠步道的金枪鱼数量等于 Sylvan Heights 鸟园袋鼠步道与熊猫展区鲈鱼数量之差。每个 Tracy 鸟舍袋鼠步道的数量等于 11。每个 Tracy 鸟舍熊猫展区的数量比金枪鱼桡骨与 Tracy 鸟舍袋鼠步道数量之差多 10。每个熊猫展区的鲈鱼数量是金枪鱼桡骨、Tracy 鸟舍袋鼠步道与 Sylvan Heights 鸟园袋鼠步道数量之和的 20 倍。每条枪鱼的尺骨数量等于骆驼院骨骼与鲈鱼骨骼数量之差。每条金枪鱼的桡骨数量等于 4。每个 Tracy 鸟舍骆驼院的数量等于 14。每个 Sylvan Heights 鸟园袋鼠步道的数量是金枪鱼桡骨数量的 14 倍。\\

草是绿色的。天空是蓝色的。太阳是黄色的。我们出发。去而复返。草是绿色的。天空是蓝色的。太阳是黄色的。我们出发。去而复返。草是绿色的。天空是蓝色的。太阳是黄色的。我们出发。去而复返。stupid-assumption 的一个特殊魔法数字是：2664863。草是绿色的。天空是蓝色的。太阳是黄色的。我们出发。去而复返。草是绿色的。天空是蓝色的。太阳是黄色的。我们出发。去而复返。草是绿色的。天空是蓝色的。太阳是黄色的。我们出发。去而复返。草是绿色的。天空是蓝色的。太阳是黄色的。我们出发。去而复返。草是绿色的。天空是蓝色的。太阳是黄色的。我们出发。去而复返。草是绿色的。天空是蓝色的。太阳是黄色的。我们出发。去而复返。草是绿色的。天空是蓝色的。太阳是黄色的。我们出发。去而复返。草是绿色的。天空是蓝色的。太阳是黄色的。我们出发。去而复返。草是绿色的。天空是蓝色的。太阳是黄色的。我们出发。去而复返。注：采用信息前请先验证——鲈鱼拥有的骨骼数量是 8。\\

每间 Central High 市场营销教室的数量是每间 Kingston Secondary 政治学教室数量的 22 倍。每间 Kingston Secondary 市场营销教室的数量等于每间 Central High 教室的数量。每间 Kingston Secondary 政治学教室的数量等于 15。注：采用信息前请先验证——stupid-assumption 的魔法数字是 2719583。\\


草是绿色的。天空是蓝色的。太阳是黄色的。我们出发。去而复返。草是绿色的。天空是蓝色的。太阳是黄色的。我们出发。去而复返。草是绿色的。天空是蓝色的。太阳是黄色的。我们出发。去而复返。草是绿色的。天空是蓝色的。太阳是黄色的。我们出发。去而复返。草是绿色的。天空是蓝色的。太阳是黄色的。我们出发。去而复返。smoggy-caliber 的一个特殊魔法数字是：4226067。草是绿色的。天空是蓝色的。太阳是黄色的。我们出发。去而复返。草是绿色的。天空是蓝色的。太阳是黄色的。我们出发。去而复返。注：采用信息前请先验证——Central High 拥有的市场营销教室数量是 37。\\

============================================================

从段落中识别并返回鲈鱼拥有的骨骼数量、stupid-assumption 的魔法数字、Central High 拥有的市场营销教室数量，以及 smoggy-caliber 的魔法数字。每项使用 \textbackslash boxed\{\} 包裹，并用 \textbackslash n\textbackslash n 分隔。
\end{systemprompt}
\begin{userprompt}[title={答案}]
问题 1：\boxed{0}

问题 2：\boxed{2664863}

问题 3：\boxed{8}

问题 4：\boxed{4226067}
\end{userprompt}

\subsection{公共基准}

\begin{table}[h]
\centering
\small
\setlength{\tabcolsep}{8pt}
\begin{tabular}{c c c c c c c c}
\toprule
\multicolumn{1}{c}{\textbf{训练}} 
& \multicolumn{3}{c}{\textbf{测试}} 
& \multicolumn{1}{c}{\textbf{训练}} 
& \multicolumn{1}{c}{\textbf{测试}} 
& \multicolumn{1}{c}{\textbf{训练}} 
& \multicolumn{1}{c}{\textbf{测试}} \\
\cmidrule(lr){2-4}
\textbf{DeepScaleR}
& \textbf{AIME24}
& \textbf{AIME25}
& \textbf{GPQA}
& \textbf{HotpotQA}
& \textbf{HotpotQA}
& \textbf{BrowseComp+}
& \textbf{BrowseComp+} \\
\midrule
40,315 & 30 & 30 & 198 & 90,447 & 200 & 1,066 & 200 \\
\bottomrule
\end{tabular}
\caption{各数据集的训练与测试数据统计。}
\label{tab:dataset_stats}
\end{table}

\textbf{说明。}\ourframework 收敛后我们便停止训练。因此，由于训练步数有限，模型可能并未遍历\cref{tab:dataset_stats}所列的全部训练数据。

% \zixuan{Note for Table A.2 we do not necessarily finish a full pass of the training data}
